Um único parágrafo que sintetize todo o meu trabalho. Organização de informação usando classificação facetada é útil para melhorar a indexação de documentos e a construção de sistemas de recuperação de informação corporativa. Essa hipótese baseia-se na evidência de facetas comuns a documentos de diferentes empresas e na flexibilidade da organização facetada. Entretanto, a classificação e indexação automáticas de um grande volume de documentos representam importantes obstáculos e nossa principal motivação. A pesquisa é descritiva, aplicada e experimental e tenta responder sobre a existência de características comuns a documentos do domínio corporativo e a possibilidade de indexação facetada automática. Duas coleções são usadas para avaliação, uma pública e outra particular. Os termos usados por autores de documentos foram obtidos através de documentos e expressões de busca. Foi empreendida uma análise preliminar do domínio corporativo, pela qual foram descobertas 12 categorias comuns e facetas úteis para o contexto de cada coleção de avaliação. A distribuição de assuntos em categorias apresentou alta correlação positiva usando o coeficiente de correlação de Spearman. Dez expressões de busca de usuários foram avaliadas no contexto da coleção particular e validaram as 12 categorias comuns. A avaliação empírica da trilha \textit{Enterprise} da \textit{Text Retrieval Conference} foi executada e os métodos de indexação, classificação e recuperação automáticos de informação facetada melhoraram a eficiência da recuperação sem fazer uso de serviços externos, como Wikipedia e metabuscadores, e sem fazer uso de estruturas hipertextuais presentes nos documentos da amostra. A avaliação empírica utilizou-se principalmente das características espaciais, temporais, de documento e de pessoal. A técnica de análise facetada mostrou-se promissora para os métodos de análise e comparação de coleções corporativas sem que dados puros sejam expostos a terceiros. A tese aponta direções de pesquisa para o uso dos métodos em outras coleções, para aperfeiçoamentos da organização da informação facetada, e para novas aplicações dos métodos também em outros domínios.